                                        IMC 2020 Online
                                       Day 2, July 27, 2020




Problem 5. Find all twice continuously dierentiable functions f : R → (0, +∞) satisfying

                                           f 00 (x)f (x) ≥ 2(f 0 (x))2

for all x ∈ R.
          Karen Keryan, Yerevan State University & American University of Armenia, Yerevan
Solution. We shall show that only positive constant functions satisfy the condition.
                  1
   Let g(x) =         . Notice that
                f (x)
                                00  0 0                    0
                          00     1    −f      2(f 0 )2 − f 00 f
                         g =        =      =                       ≤ 0,
                                 f    f2             f3

so the positive function g(x) is concave. We show that g must be constant.
    Take two arbitrary real numbers a < b. By the concavity of g , for all u < a and v > b we
have
                          g(a) − g(u)   g(b) − g(a)   g(v) − g(b)
                                      ≥             ≥             .
                             a−u           b−a           v−b
Combining this with g(u), g(v) > 0 we get
                                      g(a)   g(b) − g(a)   −g(b)
                                           >             >
                                      a−u       b−a        v−b
Now by taking limits u → −∞ and v → ∞ we obtain
                                                 g(b) − g(a)
                                            0≥               ≥ 0,
                                                    b−a
so g(a) = g(b). This holds for any pair (a, b), so g(x) is constant and f (x) = 1/g(x) also is
constant.
   If f is constant then f 0 = f 00 = 0, so the condition is satised.
   Remark. Instead of the function 1/f (x), the same idea works with arctan f (x):

                        f 00 (1 + f 2 ) − 2(f 0 )2   f 00 (1 + f 2 ) − 2(f 0 )2 (1 + f 2 )   f 00 − 2(f 0 )2
   (arctan f (x))00 =                              =                                       =                 ≥ 0.
                                (1 + f 2 )2                       (1 + f 2 )2                    1 + f2

As can be seen, arctan f (x) is a bounded convex function, therefore it must be constant.

Problem 6. Find all prime numbers p for which there exists a unique a ∈ {1, 2, . . . , p} such
that a3 − 3a + 1 is divisible by p.
                                                         Géza Kós, Loránd Eötvös University, Budapest
Solution 1. We show that p = 3 the only prime that satses the condition.




                                                        1
   Let f (x) = x3 − 3x + 1. As preparation, let's compute the roots of f (x). By Cardano's
formula, it can be seen that the roots are
     v      s 
     u
                 2    3      r                                            
     u
     3 −1     −1      −3       3    2π         2π        2π      4π       8π
 2Re      +        −      = 2Re cos    + i sin    = 2 cos , 2 cos , 2 cos
     t
        2      2      3              3          3         9       9        9

where all three values of the complex cubic root were taken.
   Notice that, by the trigonometric identity 2 cos 2t = (2 cos t)2 − 2, the map ϕ(x) = x2 − 2
cyclically permutes the three roots. We will use this map to nd another root of f , when it is
considered over Fp .
   Suppose that f (a) = 0 for some a ∈ Fp and consider
                                f (x)   f (x) − f (a)
                       g(x) =         =               = x2 + ax + (a2 − 3).
                                x−a         x−a
We claim that b = a2 − 2 is a root of g(x). Indeed,
                  g(b) = (a2 − 2)2 + a(a2 − 2) + (a2 − 3) = (a + 1) · f (a) = 0.

By Vieta's formulas, the other root of g(x) is c = −a − b = −a2 − a + 2.
   If f has a single root then the three roots must coincide, so
                                    a = a2 − 2 = −a2 − a + 2.

Here the quadratic equation a = a2 − 2, or equivalently (a + 1)(a − 2) = 0, has two solutions,
a = −1 and a = 2. By f (−1) = f (2) = 3, in both cases we have 0 = f (a) = 3, so the only
choice is p = 3.
    Finally, for p = 3 we have f (1) = −1, f (2) = 3 and f (3) = 19, from these values only f (2)
is divisible by 3, so p = 3 satises the condition.
Solution 2 (outline) Dene f (x) and g(x) like in Solution 1. The discriminant of g(x) is

                                 ∆g = a2 − 4(a2 − 3) = 12 − 3a2 .

We show that ∆g has a square root in Fp .
  Take two integers k, m (to be determinated later) and consider
         ∆g = ∆g + (ka + m)f (a) = ka4 + ma3 − (3k + 1)a2 + (k − 3m)a + (m + 12).

Our goal is to choose k, m in such a way that the last expression is a complete square. Either
by direct calculations or guessing, we can nd that k = m = 4 works:
           ∆g = ∆g + (4a + 4)f (a) = 4a4 + 4a3 − 15a2 − 8a + 16 = (2a2 + a − 4)2 .

    If p 6= 2 then we can conclude that f (x) has either no or three roots, therefore p is suitable
if and only is f (x) is a complete cube: x3 − 3x + 1 = (x − a)3 . From Vieta's formulas a3 = 1,
so a 6= 0 and 3a = 0, which is possible if p = 3.
    For p = 3 we have f (x) = (x + 1)3 , so p = 3 is suitable.
    The case p = 2 must be checked separately because the quadratic formula contains a division
by 2. f (1) = −1 and f (2) = 3, so p = 2 is not suitable.
Solution 3 (outline) Assume p > 3; the cases p = 2 and p = 3 will be checked separately.

                                                 2
   Let f (x) = x3 − 3x + 1 and suppose that a ∈ Fp is a root of f (x), and let b, c ∈ Fp2 be the
other two roots. The discriminant ∆f of f (x) can be expressed by the elementary symmetric
polynomials of a, b, c; it can be calculated that
                            ∆f = (b − c)2 (a − b)2 (a − c)2 = 81 = 92 ,

so
                                  (b − c)(a − b)(a − c) = ±9 ∈ Fp .
    Notice that ∆f 6= 0, so the three roots are distinct.
    Either b, c ∈ Fp or b, c are conjugate elements in Fp2 \ Fp , we have (a − b)(a − c) ∈ Fp , so
b − c = (b−c)(a−b)(a−c)
            (a−b)(a−c)
                         ∈ Fp . From Vieta's formulas we have b + c ∈ Fp as well; since p 6= 2, it
follows that b, c ∈ Fp . Now f (x) has three distinct roots in Fp , so p cannot be suitable.
    p = 2 does not satises the condition because both f (1) = −1 and f (2) = 3 are odd. p = 3
is suitable, because f (2) = 3 is divisible by 3 while f (1) = −1 and f (3) = 19 are not.
Problem 7. Let G be a group and n ≥ 2 be an integer. Let H1 and H2 be two subgroups of
G that satisfy

                  [G : H1 ] = [G : H2 ] = n and [G : (H1 ∩ H2 )] = n(n − 1).

Prove that H1 and H2 are conjugate in G.
   (Here [G : H] denotes the index of the subgroup H , i.e. the number of distinct left cosets
xH of H in G. The subgroups H1 and H2 are conjugate if there exists an element g ∈ G such
that g −1 H1 g = H2 .)
        Ilya Bogdanov and Alexander Matushkin, Moscow Institute of Physics and Technology
Solution 1. Denote K = H1 ∩ H2 . Since

                      n(n − 1) = [G : K] = [G : H1 ][H1 : K] = n[H1 : K],

we obtain thatF[H1 : K] = n − 1. Thus, the subgroup H1 is partitioned into n − 1 left cosets of
               i=1 hi K . Therefore, the set H1 H2 = {ab : a ∈ H1 , b ∈ H2 } is partitioned as
K , say H1 = n−1
                                    n−1
                                                 !          n−1              n−1
                                    G                       G                G
                        H1 H2 =           hi K       H2 =         hi KH2 =         hi H2 .
                                    i=1                     i=1              i=1

   The last equality holds because K ⊆ H2 , so KH2 = H2 . The last expression is a disjoint
union since
             hi H2 ∩ hj H2 6= ∅ ⇐⇒ hi−1 hj ∈ H2 ⇐⇒ h−1
                                                    i hj ∈ K ⇐⇒ hi = hj .

    Thus, H1 H2 is a disjoint union of n − 1 left cosets with respect to H2 ; hence L = G \ (H1 H2 )
is the remaining such left coset. Similarly, L is a right coset with respect to H1 . Therefore, for
each g ∈ L we have L = gH2 = H1 g , which yields H2 = g −1 H1 g .
Solution 2. Put G/H1 = X and G/H2 = Y , those are n-element sets acted on by G from the
left. Let G act on X × Y from the left coordinate-wise, consider this product as a table, with
rows being copies of X and columns being copies of Y .
    The stabilizer of a point (x, y) in X × Y is H1 ∩ H2 . By the orbit-stabilizer theorem, we
obtain that the orbit Z of (x, y) has size [G : H1 ∩ H2 ] = n(n − 1).
    If Z contains a whole column then there is a subgroup G1 of G that stabilizes x and acts
transitively on Y . If we conjugate G1 to a group G01 , then its action remains transitive on Y ,

                                                       3
so by conjugation we obtain columns of the table. Since G acts transitively on X , we cover all
the columns. It follows that Z = X × Y , so
                                n(n − 1) = |Z| = |X × Y | = n2 ,

which is a contradiction.
    Hence every column of X × Y has an element not from Z . The same holds for the rows of
X × Y . There are n elements not from Z in total and they induce a bijection between X and
Y which allows us to identify X = Y .
    After this identication, every element (x, x) from the diagonal of X × X (i.e. from (X ×
X) \ Z ) is moved to a diagonal element by any g ∈ G, because gx = gx. In this formula the
action of g in the left hand side and the action of g in the right hand side are the actions of g
on X and Y respectively.
    Therefore our bijection between X and Y is an isomorphism of sets with a left action of G.
Since H1 and H2 are stabilizers of the points in the same transitive action of G, we conclude
that they are conjugate.
    Remark. The situation in the problem is possible for every n ≥ 2: let G = Sn and let H1
an H2 be the stabilizer subgroups of two elements.




                                               4
Problem 8. Compute
                                                    n        
                                             1     X
                                                           k n
                                      lim              (−1)     log k.
                                     n→∞ log log n           k
                                                   k=1

(Here log denotes the natural logarithm.)
                                                              Fedor Petrov, St. Petersburg State University
Solution 1. Answer: 1.
   The idea is that if f (k) = g k , then
                                    R

                                                         Z
                                     X            n
                                                  k
                                            (−1)    f (k) = (1 − g)n .
                                                  k

To relate this to logarithm, we may use the Frullani integrals

    Z ∞                        Z ∞ −x        Z ∞ −kx           Z ∞ −x       Z ∞ −x
          e−x − e−kx               e              e                e            e
                     dx = lim         dx −            dx = lim         dx −        dx =
     0        x           c→+0 c    x         c     x      c→+0 c   x        kc  x
                          Z kc −x                    Z kc −x
                              e                          e −1
                      lim         dx = log k + lim             dx = log k.
                     c→+0 c     x               c→+0 c      x
   This gives the integral representation of our sum:
                             n                  Z ∞
                                                    −e−x + 1 − (1 − e−x )n
                                     
                             X       n  k
                      A :=     (−1)     log k =                            dx.
                           k=1
                                     k           0           x
   Now the problem is reduced to a rather standard integral asymptotics.
   We have (1−e−x )n ⩾ 1−ne−x by Bernoulli inequality, thus 0 ⩽ −e−x +1−(1−e−x )n ⩽ ne−x ,
and we get
          Z ∞                                           Z ∞                    Z ∞
                −e−x + 1 − (1 − e−x )n                        e−x
     0⩽                                dx ⩽ n                     dx ⩽ nM −1           e−x dx = nM −1 e−M .
            M            x                               M     x                   M

So choosing M such that M eM = n (such M exists and goes to ∞ with n) we get
                                                 Z M
                                                       −e−x + 1 − (1 − e−x )n
                                A = O(1) +                                    dx.
                                                  0             x

Note that for 0 ⩽ x ⩽ M we have e−x ⩾ e−M = M/n, and (1 − e−x )n−1 ⩽ e−e
                                                                                                      −x (n−1)
                                                                                                                 ⩽
e−M (n−1)/n tends to 0 uniformly in x. Therefore
                Z M                                                          Z M
                          (1 − e−x )(1 − (1 − e−x )n−1 )                               1 − e−x
                                                         dx = (1 + o(1))                       dx.
                  0                     x                                      0          x
Finally
                  Z M                       Z 1                  Z M               Z M
                           1 − e−x                1 − e−x              −e−x                dx
                                   dx =                   dx +              dx +              =
                      0       x              0       x            1     x              1    x
                      log M + O(1) = log(M + log M ) + O(1) = log log n + O(1),
   and we get A = (1 + o(1)) log log n.
Solution 2. We start with a known identity (a nite dierence of 1/x).


                                                          5
     Expand the rational function
                                                   m!
                                                   f (x) =
                                          x(x + 1) . . . (x + m)
as the linear combination of simple fractions f (x) = m    j=0 cj /(x + j). To nd cj we use
                                                         P

                                                                  
                                                                j m
                            cj = ((x + j)f (x)) |x=−j = (−1)           .
                                                                    j
So we get
                                       m            
                                       X
                                                   km  1            m!
                                              (−1)        =                        .                                         (1)
                                          k=0
                                                    k x+k   x(x + 1) . . . (x + m)
     Another known identity we use is
             n               n                                   
             X          n     X
                                k     k    n−1   n−1             j+1 n − 1
                  (−1)     =      (−1)         +           = (−1)          .                                                 (2)
             k=j+1
                        j    k=j+1
                                            k     k − 1                j

     Finally we write log k = 1k dx                          j=1 Ij , where Ij =              .
                                            R              Pk−1                      R 1 dx
                                 x
                                    =                                                 0 x+j
     Now we have

      n                  n           k−1
                                     X        n−1      n          X  n−1                
               n                   k n                           k n                j+1 n − 1
      X                   X                     X       X
                   k
S :=     (−1)     log k =     (−1)         Ij =     Ij      (−1)     =       Ij (−1)           =
     k=1
               k          k=1
                                     k j=1      j=1    k=j+1
                                                                   k (2) j=1              j

                n−1
             Z 1X                          Z 1        n−1                      !
                              n − 1 dx
                                    j+1             1 X         j n−1      dx
                   (−1)                    =         −      (−1)                     dx =
             0 j=1              j    x+j       0    x j=0            j    x+j           (1)

Z 1                                         Z 1                                                
      1              (n − 1)!                     dx                        1
        −                               dx =           1−                                           .
 0    x x(x + 1) . . . (x + (n − 1))           0   x       (1 + x)(1 + x/2) . . . (1 + x/(n − 1))
    So S is again expressed as an integral, for which it is not hard to get an asymptotics.
    Since et ⩾ 1+t for all real t (by convexity or any other reason), we have ey −y ⩾ 1+y 2 −y =
                                                                                   2

1+y 3
1+y
         1
      ⩾ 1+y  and 1+y
                   1
                      ⩾ e1y = e−y for y > 0. Therefore
                                                       2         1
                                                      ey −y ⩾       ⩾ e−y , y > 0.
                                                                1+y
     Using this double inequality we get
                                                                             1
                           
  x2 1+ 12 +...+     1
                                −x(1+ 12 +...+ n−1
                                                1
                                                   )                                                        1        1
 e       2         (n−1)2                              ⩾                                           ⩾ e−x(1+ 2 +...+ n−1 ) .
                                                            (1 + x)(1 + x/2) . . . (1 + x/(n − 1))
Since x2 (1 + 1/22 + . . .) ⩽ 2x2 ⩽ 2x, we conclude that
                                                                       n−1          n−1
                          1                                                1            1
                                                         , where − 2 +
                                                                       X            X
                                                   −Cn x
                                                =e                           ⩽ Cn ⩽       ,
         (1 + x)(1 + x/2) . . . (1 + x/(n − 1))                        j=1
                                                                           j        j=1
                                                                                        j

i.e., Cn = log n + O(1). Thus
             Z 1                                 Z Cn                        Z Cn          Z 1             Z Cn
                   dx                                      dt                       dt                dt                dt
       S=             (1 − e−Cn x ) =                         (1 − e−t ) =             +          −t
                                                                                              (1 − e ) +          e−t
              0    x                              0         t                 1      t      0         t     1            t
                                                = log Cn + O(1) = log log n + O(1).


                                                                      6
